 \documentclass{report}
 
 \input{~/dev/latex/template/preamble.tex}
 \input{~/dev/latex/template/macros.tex}
 
 \title{\Huge{}}
 \author{\huge{Nathan Warner}}
 \date{\huge{}}
 \fancyhf{}
 \rhead{}
 \fancyhead[R]{\itshape Warner} % Left header: Section name
 \fancyhead[L]{\itshape\leftmark}  % Right header: Page number
 \cfoot{\thepage}
 \renewcommand{\headrulewidth}{0pt} % Optional: Removes the header line
 %\pagestyle{fancy}
 %\fancyhf{}
 %\lhead{Warner \thepage}
 %\rhead{}
 % \lhead{\leftmark}
 %\cfoot{\thepage}
 %\setborder
 % \usepackage[default]{sourcecodepro}
 % \usepackage[T1]{fontenc}
 
 % Change the title
 \hypersetup{
     pdftitle={}
 }

 \geometry{
  left=1.5in,
  right=1.5in,
  top=1in,
  bottom=1in
}
 
 \begin{document}
     % \maketitle
     %     \begin{titlepage}
     %    \begin{center}
     %        \vspace*{1cm}
     % 
     %        \textbf{}
     % 
     %        \vspace{0.5cm}
     %         
     %             
     %        \vspace{1.5cm}
     % 
     %        \textbf{Nathan Warner}
     % 
     %        \vfill
     %             
     %             
     %        \vspace{0.8cm}
     %      
     %        \includegraphics[width=0.4\textwidth]{~/niu/seal.png}
     %             
     %        Computer Science \\
     %        Northern Illinois University\\
     %        United States\\
     %        
     %             
     %    \end{center}
     % \end{titlepage}
     % \tableofcontents
    \pagebreak \bigbreak \noindent
    Nate Warner \ \quad \quad \quad \quad \quad \quad \quad \quad \quad \quad \quad \quad  STAT 300 \quad  \quad \quad \quad \quad \quad \quad \quad \quad \ \ \quad Summer 2024
    \begin{center}
        \textbf{PSET 1 - Due: Sunday, June 23}
    \end{center}
    \bigbreak \noindent 
    \begin{mdframed}
        1. A small survey was conducted in which each respondent was asked how many times,
        in the previous two-week period, they had eaten at a fast food restaurant. The data appear below
        \bigbreak \noindent 
        \begin{center}
            0, 2, 1, 5, 2, 2, 3, 4, 1, 2, 7, 1, 3, 4, 1, 0, 1, 4, 2, 1, 3, 3, 2, 1, 9, 1
        \end{center}
        \begin{enumerate}[label=(\alph*)]
            \item Construct a frequency histogram. The histogram should be neat, accurate, and well labeled. 
            \item How would you describe the shape the distribution? 
            \item Find the proportion of the respondents described by each of the following. 
                \begin{enumerate}[label=(\roman*)]
                    \item Ate at a fast food restaurant at least four times
                    \item Ate at a fast food restaurant fewer than two times
                \end{enumerate}
        \end{enumerate}
    \end{mdframed}
    \bigbreak \noindent 
    Given the above data set, we can first find the frequencies of each data point. We can then construct a frequency distribution to get a better look at the data
    \bigbreak \noindent 
    \begin{center}
        \begin{tabular}{c|c}
            Number of days &  Frequency \\
            \hline
            0 & 2 \\
            1 & 8\\
            2 & 6 \\
            3 & 4\\
            4 & 3\\
            5 & 1\\
            6 & 0\\
            7 & 1\\
            8 & 0\\ 
            9 & 1
        \end{tabular}
    \end{center}
    \bigbreak \noindent 
    \tc{Number of times $n=26$ people ate a fast food restaurant in the previous two weeks organized as a frequency distribution}
    \pagebreak \bigbreak \noindent 
    From this, we create the histogram
    \bigbreak \noindent 
\begin{figure}[ht]
    \centering
    \incfig{histo22}
    \label{fig:histo22}
\end{figure}
    \fc{Number of times $n=26$ people ate a fast food restaurant in the previous two weeks organized as a histogram for discrete data}

\bigbreak \noindent 
b.) From the histogram, we can see that the shape of the distribution is skewed positively (right)
\bigbreak \noindent 
c.) Furthermore, we can find the relative frequency (proportion) of respondents who ate at a fast food restaurant at least four times by dividing the summation of the frequencies $x\in\{4,5,6,...,9\} $ by $n=26 $ (the total number of respondents). We find
\begin{align*}
    \frac{3 + 1 + 0 + 1 + 0 + 1}{26}  = 0.2308 \approx  23\%
.\end{align*}
\bigbreak \noindent 
We preform a similar calculation to find the number of respondents who ate at a fast food restaurant fewer than two times, in this case we get
\begin{align*}
    \frac{2 + 8}{26} =  0.3846 \approx 38\%
.\end{align*}






    \pagebreak 
    \begin{mdframed}
        2. The blood glucose levels (in milligrams per deciliter) for 25 patients at a medical facility appear below.
        \bigbreak \noindent 
        \begin{center}
            &63 65 66 68 69 71 73 74 75 75 76 76 77 \\
            &79 79 81 81 81 83 84 86 87 90 91 95
        \end{center}
        \bigbreak \noindent 
        \begin{enumerate}[label=(\alph*)]
            \item  Construct a relative frequency histogram. Use nine class intervals starting with 60 $\leq$ Glucose < 64. 
            \item How would you describe the shape of the distribution?
            \item What proportion of the patients had a blood glucose level described by each of the following? 
                \begin{enumerate}[label=(\roman*)]
                    \item At most 71
                    \item  Between 72 and 91 (inclusive of the endpoints)
                \end{enumerate}
        \end{enumerate}
    \end{mdframed}
    \bigbreak \noindent 
    We again begin by organizing the data via a frequency distribution
    \bigbreak \noindent 
    \begin{center}
        \begin{tabular}{c|c|c}
            Blood glucose levels $(mg/dL)$ & Frequecy & Relative frequency \\
            \hline
            60-64 & 1 &0.04\\
            64-68 & 2 &0.08\\
            68-72 & 3 &0.12\\
            72-76 & 4 &0.16\\
            76-80 & 5 &0.2\\
            80-84 & 4 &0.16\\
            84-88 & 3 &0.12\\
            88-92 & 2 &0.08\\
            92-96 & 1 & 0.04
        \end{tabular}
    \end{center}
    \tc{Blood glucose levels for $n=25$ patiens at a medical facility organized as a frequency distribution}
    \pagebreak \bigbreak \noindent 
    Using the relative frequencies found above, we construct a simple continuous data, equal class width histogram 
    \bigbreak \noindent 
\begin{figure}[ht]
    \centering
    \incfig{histo23}
    \label{fig:histo23}
\end{figure}
\bigbreak \noindent 
\fc{Blood glucose levels for $n=25$ patiens at a medical facility organized as a continous data histogram with equal class widths}
\bigbreak \noindent 
\nt{If a data point falls on the upper boundary of a class, it is instead placed in the class immediately to the right. For example 68 would be placed in the class 68-72 as opposed to 64-68}

\bigbreak \noindent 
b.) The shape of this distribution is approximately normal (symmetric)
\bigbreak \noindent 
c.) We find the proportion of patients with a glucose level of at most 71 by summing the relative frequencies from all classes in the range 60-71. This gives
\begin{align*}
    0.04 + 0.08 + 0.12 = 0.24 \approx 24\%
.\end{align*}
\bigbreak \noindent 
We do a similar calculator to find the proportion of patients with a glucose level from 72 and 91 inclusive
\begin{align*}
    0.16 + 0.2 + 0.16 + 0.12 + 0.08 = 0.72 \approx 72\%
.\end{align*}
\bigbreak \noindent 
\nt{The class 68-72 is not included, any patient with a blood glucose level of 72 would be instead placed in the class 72-76}


    \pagebreak 
    \begin{mdframed}
        3. Refer to Problem 1.
        \begin{enumerate}[label=(\alph*)]
            \item  Calculate the sample mean. 
            \item Calculate the sample median. 
            \item How do the answers to (a) and (b) compare? 
                \begin{enumerate}[label=(\roman*)]
                    \item Why should we have been able to anticipate this?  
                \end{enumerate}
        \end{enumerate}
    \end{mdframed}
    \bigbreak \noindent 
    a.) The sample mean is given by
    \begin{align*}
        \bar{x} = \frac{1}{n}&\sum x_{i} \\
        =  \frac{1}{25}&(63+65+66+68+69+71+73+74 \\
                       &+75+75+76+76+77+79+79+81 \\
                       &+81+81+83+84+86+87+90+91+95) \\
        = 77.&84
    .\end{align*}
    \bigbreak \noindent 
    \begin{remark}
        b.) The sample median is given by 
        \begin{align*}
            \tilde{x}=
            \begin{cases}
                \left(\frac{n+1}{2}\right)^{th}\ \text{ordered element} &\text{if}\ n = 2k+1,\ k\in\mathbb{Z}   \\
                \left(ave\left(\frac{n}{2}, \frac{n+2}{2}\right)\right)^{th}\ \text{ordered element} & \text{if } n=2k\, k\in\mathbb{Z} 
            \end{cases}
        .\end{align*}
    \end{remark}
    \bigbreak \noindent 
     In our case, since $n$ is odd we use case one and the median is given by
    \begin{align*}
        \tilde{x} &= \left(\frac{26}{2}\right) = 13^{th}\ \text{ordered element}  \\
        &=77
    .\end{align*}
    \bigbreak \noindent 
    c.) In this case, the sample mean $\bar{x}$ and the sample median $\tilde{x}$ are (approximately) the same. This can be anticipated when the distribution of data is symmetric (normal).
    

    \pagebreak 
    \begin{mdframed}
        4. A new type of outdoor paint was tested by painting six homes in the same geographic
        area. The number of months the paint lasted before fading was recorded and yielded the values
        below.
        \bigbreak \noindent 
        \begin{center}
            10, 60, 50, 30, 40, 20
        \end{center}
        \begin{enumerate}[label=(\alph*)]
            \item Calculate the sample range. 
            \item Calculate the sample mean. 
            \item Calculate the sample variance using the definition formula (i.e. by using the squared deviations). 
            \item Calculate the sample variance using the short-cut formula. 
            \item Calculate the sample standard deviation. 
            \item Give the units of measurement for (b); for (c); and for (e). 
        \end{enumerate}
    \end{mdframed}
    \bigbreak \noindent 
    a.) The sample range is computed by finding the difference between the largest and smallest data point. Thus, 
    \begin{align*}
        \text{Range} = 60 - 10 = 50
    .\end{align*}
    \bigbreak \noindent 
    b.) The sample mean is given by
    \begin{align*}
        \bar{x} &= \frac{1}{n}\sum x_{i} = \frac{1}{6}(10+60+50+30+40+20) \\
        &=35
    .\end{align*}
    \bigbreak \noindent 
    c.) The sample variance is given by
    \begin{align*}
        s^{2} &= \frac{1}{n-1}\sum(x_{i} -\bar{x})^{2} \\
              &= \frac{1}{5}\left((10-35)^{2}+(60-35)^{2} +(50-35)^{2}+(30-35)^{2}+(40-35)^{2}+(20-35)^{2}\right) \\
              &=350
    .\end{align*}
    \bigbreak \noindent 
    d.)  The alternative variance formula is derived by simplifying the numerator to
    \begin{align*}
        S_{xx} &= \sum x_{i}^{2} - \frac{\left(\sum x_{i}\right)^{2}}{n} \\
        &=\frac{n\sum x_{i}^{2}- \left(\sum x_{i}\right)^{2}}{n}
    .\end{align*}
    This allows us to express variance as
    \begin{align*}
        s^{2} &=\frac{1}{n-1}S_{xx} \\ 
        &=\frac{n\sum x_{i}^{2}- \left(\sum x_{i}\right)^{2}}{n^{2}-n}
    .\end{align*}
    \bigbreak \noindent 
    Using this, we again find
    \begin{align*}
        s^{2} &= \frac{6(10^{2} + 60^{2} + 50^{2} + 40^{2} + 30^{2} + 20^{2}) - (10 + 60 + 50 + 30 +40 + 20)^{2}}{6^{2}-6} \\
        &=350
    .\end{align*}
    \bigbreak \noindent 
    e.) Using the variance found above, we can then compute the std dev
    \begin{align*}
        s = \sqrt{s^{2}} = \sqrt{350} =18.7083
    .\end{align*}
    \bigbreak \noindent 
    f.) The units are months.


    \pagebreak 
    \begin{mdframed}
       5.  Refer to Problem 4.
       \begin{enumerate}[label=(\alph*)]
           \item Add 10 to each data value in Problem 4.
           \begin{enumerate}[label=(\roman*)]
               \item Re-calculate the sample mean. 
               \item How does the answer in (i) compare to that in 4(b)? Be specific. 
               \item Re-calculate the sample variance. Note – using the definition formula might be more illuminating. 
               \item How does the answer in (iii) compare to that in 4(c)? Be specific. 
           \end{enumerate}
       \item Divide each data value in Problem 4 by 10 (i.e. multiply each data value by 1/10).
           \begin{enumerate}[label=(\roman*)]
               \item  Re-calculate the sample mean. 
               \item How does the answer in (i) compare to that in 4(b)? Be specific. 
               \item Re-calculate the sample variance. Note – using the definition formula might be more illuminating. 
               \item How does the answer in (iii) compare to that in 4(c)? Be specific. 
           \end{enumerate}
       \end{enumerate}
    \end{mdframed}
    \bigbreak \noindent 
    a.i) Adding 10 to each data point gives the set
    \bigbreak \noindent 
    \begin{center}
        20 70 60 40 50 30
    \end{center}
    \bigbreak \noindent 
    Using this new set, we calculate the sample mean
    \bigbreak \noindent 
    \begin{align*}
        \bar{x} &= \frac{1}{6}(20 + 70 + 60 + 40 + 50 + 30) \\
        &=45
    .\end{align*}
    \bigbreak \noindent 
    a.ii) As expected, the new mean is exactly 10 more than that found in 4.b. The reason for this can be easily shown
    \bigbreak \noindent 
    \begin{align*}
        \bar{x} &= \frac{1}{n}\summation{n}{i=1}\ (x_{i} + 10)\ \\
                &=\frac{1}{n}\summation{n}{i=1}\ x_{i}\ + \frac{1}{n}\summation{n}{i=1}\ 10\ \\
                &=\frac{1}{n}\summation{n}{i=1}\ x_{i}\ + \frac{1}{n}10n \\
                &=10+\frac{1}{n}\summation{n}{i=1}\ x_{i}
    .\end{align*}
    \bigbreak \noindent 
    Thus, finding the mean after adding 10 to  each data point is exactly the same as adding 10 to the original mean.

    \bigbreak \noindent 
    a.iii) Using the new data set and the new sample mean, we can compute the new sample variance.
    \begin{align*}
        s_{y}^{2} &= \frac{1}{n-1}\summation{n}{i=1}\ (y_{i} - \bar{y})^{2}\ \\
                  &= \frac{1}{5}((20-45)^{2} + (70-45)^{2} + (60-45)^{2}  \\
                  &+ (40-45)^{2} + (50-45)^{2} + (30-45)^{2}) \\
                  &=350
    .\end{align*}
    \bigbreak \noindent 
    a.iv) Here we see that the variance remains unchanged. The reason for this can be shown mathematically
    \begin{align*}
        s^{2}&=\frac{1}{n-1}\summation{n}{i=1}\ ((x_{i} + 10) - (\bar{x} +10))^{2}\ \\
        &=\frac{1}{n-1}\summation{n}{i=1}\ (x_{i}+10-\bar{x}-10)^{2}\ \\
        &=\frac{1}{n-1}\summation{n}{i=1}\ (x_{i}-\bar{x})^{2}\
    .\end{align*}
    \bigbreak \noindent 
    Thus, we notice adding a constant to each value in a data set produces no change in the variance.
    \bigbreak \noindent 
    b) Multiply each value in the original data set by $\frac{1}{10}$ gives the new set
    \begin{center}
        1 6 5 3 4 2
    \end{center}
    \bigbreak \noindent 
    b.i) With this, we can calculate the new sample mean
    \begin{align*}
        \bar{x} &= \frac{1}{6}(1+6+5+3+4+2) \\
        &=\frac{7}{2} =3.5
    .\end{align*}
    \bigbreak \noindent 
    b.ii) As we see here, this new sample mean is exacty the original sample mean divided by 10. This can again be shown using simple properties of summation. In specific we know $\sum cx_{i} = c\sum x_{i}$ 
    \begin{align*}
        \bar{x} &= \frac{1}{n}\summation{n}{i=1}\ \left(x_{i}\cdot \frac{1}{10}\right)\  \\
       &=\frac{1}{10n}\summation{n}{i=1}\ x_{i}\ \\
       &= \frac{\frac{1}{n}\summation{n}{i=1}\ x_{i}\ }{10}
    .\end{align*}
    \bigbreak \noindent 
    b.ii) Thus, we see that multiply or dividing a data set by a constant factor also multiplys or divides the mean by that same constant factor.
    \bigbreak \noindent 
    b.iii) Computing the new sample variance gives
    \begin{align*}
        s^{2} &= \frac{1}{5}((1-3.5)^{2} + (6-3.5)^{2} \\
       &+ (5-3.5)^{2} + (3-3.5)^{2} + (4-3.5)^{2} + (2-3.5)^{2})\\ 
       &=3.5  \
    .\end{align*}
    \bigbreak \noindent 
    b.iv) We see that the new sample variance does actually change this time, by a factor of $\frac{1}{10^{2}}$. Again, we can show why this is the case
    \begin{align*}
        s_{2}^{2} &= \frac{1}{n-1}\summation{n}{i=1}\ \left(\frac{x_{i}}{10} - \frac{\bar{x}}{10}\right)^{2} \\
        &=\frac{1}{n-1}\summation{n}{i=1}\ \left(\frac{1}{10}(x_{i}-\bar{x})\right)^{2}\ \\
        &=\frac{1}{n-1}\summation{n}{i=1}\ \left(\frac{1}{10}\right)^{2}(x_{i}-\bar{x})^{2}\ \\
        &=\frac{1}{10^{2}}\left(\frac{1}{n-1}\summation{n}{i=1}(x_{i}-\bar{x})^{2} \right)
    .\end{align*}

    \pagebreak 
    \begin{mdframed}
        6. Suppose that $x_{1}, x_{2}, ..., x_{n}$ is a set of data with $\bar{x}= 40$, $\tilde{x} = 60$, $s_{x}^{2} = 25$ and $s_{x} = 5 $
        Suppose that each $x_{i}$
        is transformed to a new data value $y_{i} = 2x_{i} - 1$
        \begin{enumerate}[label=(\alph*)]
            \item Compare $\bar{x}$ to $\tilde{x}$. What do the values potentially tell us about the shape of the distribution? 
            \item Find each of the following
                \begin{enumerate}[label=(\roman*)]
                    \item $\bar{y}$ 
                    \item $\tilde{y}$
                    \item $s_{y}^{2}$
                    \item $s_{y}$
                \end{enumerate}
        \end{enumerate}
    \end{mdframed}
    \bigbreak \noindent 
    6.a) Since the mean and median are far apart, specifically the median being larger than the mean, we can expect the shape of the distribution to be negatively skewed
    \bigbreak \noindent 
    \begin{remark}
                    Assume that \( Y \) is a linear transformation of \( X \). Then,
            \bigbreak \noindent
            \[
                Y = bX + A
            \]
            \bigbreak \noindent 
            \begin{itemize}
                \item Mean of \( Y \): 
                    \[
                        \text{Mean of } Y = b (\text{Mean of } X) + A
                    \]
                \item Median of \( Y \): 
                    \[
                        \text{Median of } Y = b (\text{Median of } X) + A
                    \]
                \item Standard Deviation of \( Y \): 
                    \[
                        \text{sd of } Y = b (\text{sd of } X)
                    \]
                \item Variance of \( Y \): 
                    \[
                        \text{Variance of } Y = b^2 (\text{Variance of } X)
                    \]
            \end{itemize}
    \end{remark}
    \bigbreak \noindent 
    b.i) The new mean can be calculated by
    \begin{align*}
        2(40)-1 = 79
    .\end{align*}
    \bigbreak \noindent 
    b.ii) Median:
    \begin{align*}
       2(60)-1 = 119 
    .\end{align*}
    \bigbreak \noindent 
    b.iii) Variance:
    \begin{align*}
        4(25) = 100
    .\end{align*}
    \bigbreak \noindent 
    b.iv) Standard deviation:
    \begin{align*}
        2(5) = 10
    .\end{align*}


 \end{document} % (:
